% Unicode corpus (lane UNICODE): a real document rather than a feature
% probe. A short paper that loads everything at once, as authors do:
% fontspec by family name (TeX Gyre Pagella and Heros), a project font
% (JetBrains Mono, OFL, from the project's fonts/ folder by Path=),
% unicode-math (TeX Gyre Pagella Math), polyglossia with Greek (Libertinus
% Serif by name), Hebrew and Arabic (bidi), xeCJK (Fandol), hyperref, a
% table of contents, footnotes, a table, cross references and a
% bibliography. Every font is one the no-TeX-Live bundle packs (GFS Didot,
% tried first, is not: tools/bundle/tl2026/extra.txt).
\documentclass[11pt,a4paper]{article}
\usepackage{amsmath}
\usepackage{fontspec}
\usepackage{unicode-math}
\setmainfont{TeX Gyre Pagella}
\setsansfont{TeX Gyre Heros}[Scale=MatchLowercase]
\setmonofont{JetBrainsMono}[Path=fonts/, Extension=.ttf,
  UprightFont=*-Regular, BoldFont=*-Bold, ItalicFont=*-Italic,
  BoldItalicFont=*-BoldItalic, Scale=MatchLowercase]
\setmathfont{TeX Gyre Pagella Math}
\usepackage{xeCJK}
\setCJKmainfont[BoldFont=FandolSong-Bold.otf, ItalicFont=FandolKai-Regular.otf]{FandolSong-Regular.otf}
\setCJKsansfont[BoldFont=FandolHei-Bold.otf]{FandolHei-Regular.otf}
\usepackage{booktabs}
\usepackage{xcolor}
\usepackage[colorlinks=true,linkcolor=blue!60!black,citecolor=green!40!black]{hyperref}
\usepackage{polyglossia}
\setdefaultlanguage{english}
\setotherlanguage{greek}
\setotherlanguage{hebrew}
\setotherlanguage{arabic}
\newfontfamily\greekfont{Libertinus Serif}
\newfontfamily\hebrewfont{DavidCLM}[Script=Hebrew, Extension=.otf,
  UprightFont=*-Medium, BoldFont=*-Bold, ItalicFont=*-MediumItalic,
  BoldItalicFont=*-BoldItalic]
\newfontfamily\arabicfont{Amiri}[Script=Arabic, Extension=.ttf,
  UprightFont=*-Regular, BoldFont=*-Bold, ItalicFont=*-Italic,
  BoldItalicFont=*-BoldItalic]
\title{Typesetting Many Scripts in One Document}
\author{A. Author\thanks{Department of Typography, \texttt{author@example.org}.}}
\date{1 November 2026}
\begin{document}
\maketitle
\begin{abstract}
We report on a document that mixes Latin, Greek, Hebrew, Arabic and Chinese
text with OpenType mathematics. The aim is ordinary: a paper an author would
write, with a table of contents, footnotes, a table, cross references and a
short bibliography~\cite{knuth, bringhurst}.
\end{abstract}
\tableofcontents

\section{Introduction}\label{s:intro}
Real documents rarely exercise one feature at a time. This one sets its body
in Pagella, its headings in Heros (\textsf{like this, \textbf{bold}}) and its
code in a font the project ships: \texttt{fn main() \{ println!("hi"); \}},
\texttt{\textbf{bold}}, \texttt{\textit{italic}}. Ligatures (fi, fl, ffi),
old-style figures {\addfontfeatures{Numbers=OldStyle}0123456789} and small
capitals (\textsc{Small Caps}) come from the main font.\footnote{Pagella's
small capitals are real glyphs, not scaled capitals.} Section~\ref{s:math}
has the mathematics and Section~\ref{s:scripts} the other scripts.

\section{Mathematics}\label{s:math}
Inline, the Gaussian integral $\int_{-\infty}^{\infty} e^{-x^2}\,dx=\sqrt{\pi}$
and a sum $\sum_{k=1}^{n} k = \frac{n(n+1)}{2}$ sit in the line. Displayed,
with Unicode input:
\begin{equation}\label{eq:euler}
  e^{iπ} + 1 = 0, \qquad ∀ε>0\ ∃δ>0:\ |x-a|<δ ⇒ |f(x)-f(a)|<ε .
\end{equation}
\begin{align}
  \nabla \times \mathbf{B} &= μ_0 \mathbf{J} + μ_0 ε_0 \frac{∂\mathbf{E}}{∂t}, \label{eq:maxwell}\\
  \mathbb{P}(A \mid B) &= \frac{\mathbb{P}(B \mid A)\,\mathbb{P}(A)}{\mathbb{P}(B)} .
\end{align}
Equations~\eqref{eq:euler} and~\eqref{eq:maxwell} use \verb|unicode-math|'s
symbols: $\mathcal{L}$, $\mathfrak{g}$, $\mathsf{T}$, $\symbf{x}$, $ℝ^n$, $ℤ$.

\section{Other scripts}\label{s:scripts}
\subsection{Greek}
\begin{greek}
Ἐν ἀρχῇ ἦν ὁ λόγος, καὶ ὁ λόγος ἦν πρὸς τὸν θεόν.
\end{greek}

\subsection{Hebrew and Arabic}
Hebrew inline, \texthebrew{שלום עולם}, then a paragraph.

\begin{hebrew}
זהו טקסט בעברית עם מספרים 2026 ו\textbf{הדגשה}. משפט שני, ארוך יותר, כדי
שהפסקה תתפרס על פני יותר משורה אחת.
\end{hebrew}

Arabic inline, \textarabic{السلام عليكم}, then a paragraph.

\begin{Arabic}
مرحبا بالعالم. هذا نص عربي مع \textbf{خط عريض} وأرقام ١٢٣٤. هذه فقرة
أطول قليلا حتى تمتد على أكثر من سطر واحد.
\end{Arabic}

\subsection{Chinese}
中文排版：这是一个中文段落，包含标点符号“引号”、（括号）和《书名号》。中文与
English混排，数字123和公式$a^2+b^2=c^2$。\textbf{粗体}、\textit{楷体}、
\textsf{黑体}。这一段要足够长，才能检验中文的断行与标点挤压。

\section{Results}\label{s:results}
Table~\ref{t:scripts} lists the scripts and the fonts that set them.
\begin{table}[h]
\centering
\begin{tabular}{llr}
\toprule
Script & Font & Sample \\
\midrule
Latin   & TeX Gyre Pagella & Ågård \\
Greek   & Libertinus Serif & \textgreek{λόγος} \\
Hebrew  & David CLM        & \texthebrew{עברית} \\
Arabic  & Amiri            & \textarabic{عربي} \\
Chinese & Fandol Song      & 中文 \\
Code    & JetBrains Mono   & \texttt{a->b} \\
\bottomrule
\end{tabular}
\caption{Scripts and fonts.}\label{t:scripts}
\end{table}
The links in this paragraph go to Section~\ref{s:intro}, to
Equation~\eqref{eq:euler} and to \href{https://ctan.org/pkg/fontspec}{fontspec on CTAN}.

\begin{thebibliography}{9}
\bibitem{knuth} D.~E. Knuth, \emph{The \TeX book}, Addison-Wesley, 1984.
\bibitem{bringhurst} R.~Bringhurst, \emph{The Elements of Typographic Style}, Hartley \& Marks, 2004.
\end{thebibliography}
\end{document}
